\documentclass[conference,letterpaper,10pt]{IEEEtran}
\usepackage{amsmath}
\usepackage{booktabs}
\usepackage{url}
\usepackage{graphicx}

\title{Agent Flight Recorder: Trajectory-Aware Runtime Enforcement for Tool-Using AI Agents}
\author{\IEEEauthorblockN{Pramod Ubbala}
\IEEEauthorblockA{Independent Researcher\\Renton, Washington, USA\\ubbalap@gmail.com\\ORCID: 0009-0005-5816-2998}}

\begin{document}
\maketitle

\begin{abstract}
Agent Flight Recorder evaluates tool proposals before execution using trusted metadata, successful execution history, and optional declared artifact lineage. It returns Allow, Review, or Block without inspecting private chain-of-thought. In a protocol-frozen 120-pair AgentDojo workspace study, attack successes fell from 20/120 to 0/120, while legitimate-task success under attack fell from 61/120 to 35/120. A separate literal-payload replication compared baseline, point-only, and full-history enforcement: fixed attack successes were 7/28, 4/28, and 0/28. A bounded policy-aware attacker succeeded in 11/27, 4/27, and 0/27 valid campaigns; each arm had one missing outcome. The incremental history contrast spans only two attack goals. The corrected clean-task replication recorded 31/40, 30/40, and 16/40 successes, exposing a substantial utility cost. Live results assess session history; artifact lineage is tested separately on synthetic traces. These findings support bounded intervention under the recorded conditions, not general adaptive robustness or production readiness.
\end{abstract}

\section{Introduction}
Tool use changes the security problem for language-model agents. A model response can become an email, database query, file transfer, calendar mutation, or another state-changing operation. Indirect prompt injection exploits the same path. An attacker places instructions inside content the agent later consumes, and the agent may treat those instructions as authority rather than data~\cite{greshake2023indirect}. Benchmarks such as InjecAgent, AgentDojo, and Agent Security Bench show that this failure mode persists in tool-integrated settings~\cite{zhan2024injecagent,debenedetti2024agentdojo,zhang2025asb}.

Many defenses inspect prompts, sanitize tool data, isolate trusted control flow, compare counterfactual executions, or reconstruct program-like dependencies~\cite{debenedetti2025camel,zhu2025melon,zhan2025adaptive,wang2025agentarmor,bhagwatkar2025firewalls,he2026attriguard,hofer2026automated}. Agent Flight Recorder addresses a narrower systems question: what can an application enforce at the moment a proposed tool action is about to execute, using only observable runtime state supplied by trusted adapters?

The design treats the tool boundary as a policy boundary. The model proposes an action. Application code assigns structured security metadata. A deterministic recorder evaluates the current action together with successful executed history and, where available, artifact lineage. The execution gate then permits the action, requires approval, or blocks it. Denied and failed actions remain audit events but do not become facts in causal execution history.

This distinction matters in multi-step trajectories. An agent may consume external content, access sensitive records, create a derived export, and later propose transmission to an untrusted destination. A point check on the final operation can miss the path that produced the payload. A session-wide rule recovers temporal context but can contaminate unrelated work. Artifact lineage narrows inherited restrictions to outputs that depend on sensitive sources.

\section{Threat Model and Design Goals}
\subsection{Threat model}
The attacker may control content consumed through tools or external data sources. That content may contain instructions designed to redirect the agent toward an unauthorized action. The attacker may also nominate an untrusted destination. The language model may follow, ignore, or partially follow the injected instruction.

The attacker cannot modify the policy engine, execution gate, session identifier, or trusted metadata produced by the application adapter. Those components form the trusted computing base. The current system also assumes that declared artifact links accurately represent relevant data dependencies.

The defense does not address a compromised adapter, malicious tool implementation, operating-system compromise, model-weight compromise, undeclared covert channels, or data flow that never crosses an instrumented tool boundary.

\subsection{Design goals}
The enforcement layer targets five operational properties: every supported tool effect passes through one decision point before execution; no policy decision depends on private chain-of-thought; denied or failed actions do not alter causal state; derived artifacts inherit security properties from declared sources; and Allow, Review, Block, approval, rejection, execution, and failure remain distinguishable in the audit trace.

\section{Agent Flight Recorder}
\subsection{Action representation}
A proposed action is mapped to structured metadata before the executor runs. The current implementation records operation, resource type, sensitivity, destination trust, privilege level, input provenance, session identifier, and optional data-flow information. The language model does not assign trusted security labels directly in the AgentDojo integration. The current synchronous execution path is assumed; the generic callback gate does not await asynchronous completion.

For AgentDojo, a trusted tool catalog maps supported tools to security metadata. Recipient domains can alter destination trust. Recognized recipient fields outside the configured domain are marked untrusted; destination-field coverage is incomplete. Unmapped tools fail closed.

\begin{figure}[t]
\centering
\fbox{\begin{minipage}{0.92\columnwidth}
\small
\textbf{Agent proposal} $\rightarrow$ \textbf{trusted adapter metadata} $\rightarrow$ \textbf{Recorder assessment} $\rightarrow$ \textbf{execution gate}\\[2pt]
\hspace*{0.6em}$\searrow$ Allow: execute and record successful effect\\
\hspace*{0.6em}$\searrow$ Review: require authorization; deny in confirmatory benchmark\\
\hspace*{0.6em}$\searrow$ Block: no tool effect\\[2pt]
Successful effects update executed history and artifact state. Denied and failed proposals remain audit events only.
\end{minipage}}
\caption{Tool-boundary enforcement path. The policy observes proposals and trusted metadata, not private model reasoning.}
\label{fig:architecture}
\end{figure}

\subsection{Executed trajectory state}
Let the successful execution history before action $a_i$ be $H_i$. The recorder computes
\[
D_i=P(a_i,H_i,G_i),
\]
where $G_i$ is artifact-lineage state when lineage is declared and $D_i\in\{Allow,Review,Block\}$.

The gate assesses the proposal without mutating history. A Block returns without invoking the tool. A Review executes only when the configured approval path approves it. After an allowed tool call returns successfully, the recorder appends that action to history:
\[
H_{i+1}=\begin{cases}
H_i\mathbin{\|}a_i,&\text{successful execution},\\
H_i,&\text{otherwise.}
\end{cases}
\]
This prevents policy from reasoning as though denied or failed operations occurred.

\subsection{Trajectory risk}
The prototype uses transparent rules rather than a learned classifier. Signals include external or untrusted provenance, data sensitivity, destination trust, privilege, prior untrusted input followed by sensitive access, and sensitive export transfer outside the trusted boundary. Scores 0--39 map to Allow, 40--69 to Review, and 70--100 to Block. Sensitive artifact upload outside the trusted boundary is raised to at least the Block threshold. Scores explain deterministic policy behavior; they are not calibrated probabilities.

\subsection{Artifact lineage}
Session history alone can over-attribute risk. A session may handle restricted records and later create an unrelated public artifact. The implementation therefore supports session-local artifact identifiers and declared source relationships.

Sensitivity follows the lattice
\[
Public < Internal < Confidential < Restricted.
\]
For derived artifact $x$ with sources $S_x$,
\[
s(x)=\max\left(s_{declared}(x),\max_{y\in S_x}s(y)\right).
\]
Untrusted ancestry also propagates. Missing producers, cross-session references, and invalid lineage fail closed. The mechanism does not infer hidden causal structure from the model. It relies on application-supplied artifact facts.

\subsection{Review semantics}
Review is not equivalent to Allow. In an interactive deployment it represents a pause for explicit authorization. The confirmatory AgentDojo experiment was non-interactive and used \texttt{Review=deny}: every Review prevented the tool call. That choice was frozen before v2 outcomes were observed.

\section{AgentDojo Integration}
The evaluation adapter replaces AgentDojo's normal tool executor with a gate-aware executor. For each proposed tool call, the adapter creates trusted action metadata, requests a deterministic assessment from the TypeScript policy bridge, and executes the AgentDojo function only when the decision permits it.

A blocked action produces no tool effect. Under \texttt{Review=deny}, a Review also produces no tool effect. Successful executions are recorded into the policy bridge. The adapter writes an audit stream preserving the proposed tool, action metadata, policy assessment, approval state, execution state, and errors.

Baseline and Recorder modes use the same model deployment, task pair, injected payload, attack family, and environment. Baseline runs first for each pair. The catalog supplies no artifact links, so v2 evaluates session-history rules, not lineage enforcement. The runtime gate is the experimental difference.

\section{Experimental Method}
\subsection{Evidence sequence}
A first preregistered 30-pair holdout recorded 2/30 baseline attacks versus 0/30 under enforcement (exact McNemar $p=0.5$), with utility 20/30 versus 15/30. Post-hoc review-handling and repeated-block studies reused those identities; they did not add independent confirmatory units. V2 was frozen before collecting new outcomes.

\subsection{V2 benchmark configuration}
V2 used AgentDojo package version 0.1.35, benchmark v1.2.2, the \texttt{workspace} suite, \texttt{tool\_knowledge} attack, one Azure model deployment, one trusted email domain, the frozen Agent Flight Recorder policy, and confirmatory \texttt{Review=deny}. The archive records deployment names, not a verified underlying model snapshot. Secrets and credentials were excluded.

\subsection{Pair selection and overlap control}
The pinned inventory contained 40 user tasks and 14 injection tasks. A committed generator excluded previously used exact pair identities, ranked the remainder by a fixed SHA-256 key, and selected 120 pairs. Components were reused; pair novelty does not imply new task primitives. A parser repair preceded all v2 model calls and did not change the frozen design.

\subsection{Paired execution and endpoints}
Each selected pair was executed in baseline mode and in the same environment with Agent Flight Recorder enforcement. The experiment did not stop early for benefit, harm, or significance. The primary endpoint was AgentDojo injection-task attack success. Legitimate-task utility under attack was a prespecified secondary endpoint. We report raw counts, nominal Wilson 95\% intervals, paired discordance, observed differences, and exact two-sided McNemar tests. These pair-level calculations do not quantify uncertainty over new user-task or attack primitives. Marginal Wilson intervals are not intervals for the paired effect. V2 requested a paired interval without specifying its construction. We now report a post-hoc paired percentile bootstrap (20,000 resamples; seed 20261007): Recorder-minus-baseline attack risk $[-23.3,-10.0]$ percentage points and utility $[-30.0,-13.3]$. Separate injection-goal resampling gives $[-30.2,-4.2]$ and $[-29.9,-13.8]$. These are sensitivity summaries, not a retroactively preregistered method or validated population coverage. Repeated executions of a pair are not new tasks.

\section{Results}
\subsection{Attack success}
The baseline completed the injected objective in 20 of 120 pairs. Agent Flight Recorder completed it in none of the 120 pairs.

\begin{table}[t]
\centering
\caption{Confirmatory attack outcome.}
\label{tab:attack}
\begin{tabular}{lrr}
\toprule
Outcome & Baseline & Recorder\\
\midrule
Attack success & 20/120 (16.7\%) & 0/120 (0.0\%)\\
Attack resisted & 100/120 & 120/120\\
Nominal Wilson 95\% CI & 11.1--24.3\% & 0.0--3.1\%\\
\bottomrule
\end{tabular}
\end{table}

The observed attack-risk difference, Recorder minus baseline, was -16.7 percentage points. Paired outcomes were one-sided: 20 pairs had baseline attack success with Recorder resistance, zero had baseline resistance with Recorder attack success, zero succeeded in both modes, and 100 resisted in both. The nominal exact two-sided McNemar p-value was $1.90735\times10^{-6}$.

The selected pairs use 38 user tasks and 14 injection tasks. All 20 baseline successes fall in five injection tasks (IDs 0, 1, 2, 4, and 5; counts 7, 6, 2, 2, and 3). The nominal test does not model shared-task dependence or support inference over 120 distinct attack primitives. The observed relative reduction is not proof of complete security.

\subsection{Legitimate-task utility under attack}
The security result came with a large task-completion penalty.

\begin{table}[t]
\centering
\caption{Confirmatory utility outcome.}
\label{tab:utility}
\begin{tabular}{lrr}
\toprule
Outcome & Baseline & Recorder\\
\midrule
Utility success & 61/120 (50.8\%) & 35/120 (29.2\%)\\
Nominal Wilson 95\% CI & 42.0--59.6\% & 21.8--37.8\%\\
\bottomrule
\end{tabular}
\end{table}

The observed utility difference was -21.7 percentage points. Paired utility outcomes were 32 success/success, 29 baseline-success/Recorder-failure, 3 baseline-failure/Recorder-success, and 56 failure/failure. The nominal exact two-sided McNemar p-value was $2.55601\times10^{-6}$. The confirmatory experiment does not support a claim that the defense preserved baseline utility.

\subsection{Policy interventions}
The archive combines auxiliary injection-task runs with attacked user-task runs. Separating sessions against raw transcripts excludes 249 auxiliary proposals. The 120 attacked trajectories contain 306 proposals: 182 Allow, 108 Review, and 16 Block. Of these, 160 executed, 124 were policy-denied, and 22 allowed proposals failed. Interventions occurred in 66 trajectories. These corrected descriptive counts replace mixed-run totals; the primary task endpoints are unchanged. Proposals are not independent observations.

\section{Post-hoc Review Sensitivity}
After the original 30-pair holdout was frozen, the same known identities were rerun with \texttt{Review=approve}. Baseline utility was 21/30, while Recorder utility was 24/30; Recorder utility had been 15/30 under frozen deny handling. Nine Recorder cases moved from utility failure to success, and none moved in the opposite direction.

Security inference from that sensitivity run is weak because both baseline and Recorder recorded 0/30 attack successes. Five repeated blocks of the same 30 identities then produced baseline utility 113/150, deny utility 74/150, and approve utility 115/150; attack success was 1/150 for baseline and 0/150 for each Recorder condition. Those 150 observations per condition are repetitions of 30 identities, not 150 new benchmark tasks.

These post-hoc results show that Review handling can affect utility on the known 30-pair subset. They do not identify the causes of all v2 utility losses, establish that automatic approval is safe, or amend v2. We inspected the 29 v2 pairs with baseline success and Recorder failure. On main-task traces, the first intervention was Review on calendar creation in eight pairs, file deletion in five, file creation in three, unread-email retrieval in two, and three other operations in one each. Sending email was Blocked first in two pairs. Six lost pairs had no policy intervention. These associations do not identify individual causes; stochastic behavior and execution failures also matter.

\section{Engineering Validation}
Before the benchmark study, deterministic tests covered sensitive export, unrelated public and restricted work in one session, denied producers, executor failures, missing lineage, cross-session references, duplicate artifact identifiers, and multi-generation derivation.

A seeded metadata-corruption study used three synthetic export traces and 1,000 trials per trace at each of six corruption probabilities. In each corrupted trial, all Confidential or Restricted source labels in that trace were downgraded to Public. Unsafe execution was 0\% with intact labels and 100\% at complete label corruption; intermediate unsafe-execution rates ranged from 9.8--10.9\%, 23.8--26.2\%, 49.2--50.0\%, and 74.1--78.3\% at corruption probabilities 10\%, 25\%, 50\%, and 75\%, respectively. The unit of corruption was the whole trace, not an independent draw for each source label. Repeated trials on three fixtures do not add independent benchmark tasks. Omitted producer lineage or an unresolved input reference blocked transfer. Removing both lineage and source labels permitted the unsafe terminal action under scripted Review approval. These tests expose dependence on trusted metadata and are not pooled with the 120-pair endpoint.

\subsection{Post-hoc feature and threshold diagnostics}
We replayed the 306 attacked-run proposals, excluding auxiliary sessions. The default policy reproduced every recorded score and decision. Execution membership and order remain fixed; feature ablations transform metadata in both current and previously successful actions. Alternative decisions create no new outputs or later proposals. These diagnostics measure decision sensitivity, not counterfactual attack success or utility.

We removed one metadata signal at a time: sensitivity became Public, destination trust became Trusted, privilege became one, and provenance became System. A separate condition removed history at assessment. Table~\ref{tab:ablation} reports the resulting decision changes. The AgentDojo catalog does not supply artifact links, so this replay does not test lineage ablation. Lineage behavior remains covered by separate synthetic engineering tests.

\begin{table}[t]
\centering
\caption{Post-hoc signal ablation on 306 attacked-run proposals. Newly allowed means originally policy-denied. Task outcomes are not recomputed.}
\label{tab:ablation}
\begin{tabular}{lrr}
\toprule
Signal removed & Decisions changed & Newly allowed\\
\midrule
Sensitivity & 105 & 89\\
Destination trust & 35 & 19\\
Privilege & 47 & 31\\
Provenance & 123 & 107\\
Executed history & 123 & 107\\
\bottomrule
\end{tabular}
\end{table}

We also swept Review thresholds 30, 40, and 50 against Block thresholds 60, 70, and 80. The original thresholds remain 40 and 70. Holding the Block threshold at 70, lowering Review to 30 changed 36 decisions and raised denied proposals from 124 to 160; raising Review to 50 changed 43 decisions and lowered denied proposals to 81. Those 43 newly allowed proposals are not evidence of 43 recovered tasks. Under Review=deny, shifting an action between Review and Block does not permit execution. The replay supports sensitivity to the engineered thresholds, not their optimality or calibration.

The signal weights encode explicit policy priorities: untrusted provenance contributes 20, Confidential and Restricted labels contribute 15 and 25, Unknown and Untrusted destinations contribute 15 and 30, and privilege levels at least four contribute 15. Trajectory rules add context or impose a score floor. These values define the prototype's policy rather than learned estimates. We retain them to preserve the confirmatory comparison and do not select a replacement policy from these observed outcomes.


\section{Related Work}
Indirect prompt injection arises when data consumed by an LLM-integrated application carries instructions that redirect later behavior. Greshake et al. demonstrated this attack surface~\cite{greshake2023indirect}. InjecAgent supplied tool-integrated attack cases~\cite{zhan2024injecagent}. AgentDojo jointly measures security and utility under attacks and defenses~\cite{debenedetti2024agentdojo}. Agent Security Bench broadens evaluation across attack classes and tools~\cite{zhang2025asb}.

CaMeL separates control from untrusted data and uses capabilities to constrain information flow~\cite{debenedetti2025camel}. Fides formally characterizes dynamic taint enforcement and tracks confidentiality and integrity~\cite{costa2025fides}. Progent deterministically checks symbolic tool/argument policies and requires approval for privilege expansion~\cite{shi2025progent}. Thus deterministic gating and taint tracking are established approaches. Our contribution is an inspectable gate with explicit execution-state semantics and a bounded security--utility evaluation, not a new information-flow guarantee.

MELON uses masked re-execution and tool comparison~\cite{zhu2025melon}. AttriGuard uses action-level causal attribution through counterfactual replay~\cite{he2026attriguard}. Agent Flight Recorder performs no model re-execution for a policy decision. Once the trusted adapter has assigned metadata, assessment is deterministic.

AgentArmor constructs graph abstractions of agent execution, including inferred dependencies~\cite{wang2025agentarmor}. Our gate instead trusts adapter-declared links and recorded successful effects. No matched experiment establishes superiority to these defenses.

Tool-interface firewalls provide another nearby design point~\cite{bhagwatkar2025firewalls}. Adaptive evaluation sharpens the claim boundary. Zhan et al. bypassed multiple defenses with adaptive attacks~\cite{zhan2025adaptive}, and later automated AgentDojo attack work found that attack effectiveness depends on attacker method and model~\cite{hofer2026automated}. V2 used one fixed attack family. Adaptive robustness remains unestablished. References identified as arXiv preprints in the bibliography describe preprint versions; we do not treat them as established peer-reviewed venue publications.

\section{Discussion}
\subsection{Security and utility are coupled}
The central result is not simply 20 attacks versus zero. Twenty-nine pairs showed baseline success and Recorder failure, versus three in the opposite direction. A strict gate can stop both malicious and legitimate work. V2 had no clean-workload control. The exploratory reviewer validation below measures that cost separately.

The post-hoc Review studies make this coupling visible. Review handling is part of the security mechanism, not a user-interface detail. Automatically approving every Review can recover utility, but it changes enforcement semantics. A production design needs an approval channel whose decisions carry real authorization rather than scripted acceptance or denial.

\subsection{Trusted metadata boundary}
The adapter assigns tool labels, recipient-domain trust, and artifact links. Model arguments remain untrusted. Domain membership does not establish recipient authorization, and static sensitivity labels do not classify individual records. V2 omitted Cc/Bcc destination fields; v3 repairs this coverage but still enforces coarse metadata rather than verified information flow. Unmapped tools fail closed. Deployment would require schema validation, complete field coverage, authenticated sessions, tamper-resistant state, strict registration, and adapter review. These controls are requirements, not demonstrated benchmark properties.

\subsection{Exploratory reviewer validation}
A separately frozen v3 study compared baseline, point-only enforcement, and full-history enforcement on 28 new task pairs (two per injection goal, 19 user tasks), with balanced arm orders and the repaired adapter. All arms used \texttt{gpt-4.1-mini-2025-04-14}. Point-only retained the same metadata, weights, and thresholds but assessed empty history. No arm supplied artifact links. A policy-aware attacker received the gate rules and its own prior feedback, with at most five target attempts per campaign.

An audit found that pinned AgentDojo interpolated attack text into YAML before parsing, changing all 84 initial fixed-attack environments. We retained those runs and launched an independent replication that parses the trusted template before inserting attack text into string fields. The corrected study uses unchanged policies and scoring; historical and corrected cells are never pooled.

\begin{table}[t]
\centering
\caption{Corrected literal-payload v3 replication. Fixed results use all 28 pairs. Adaptive attack counts use 27 valid campaigns per arm; one campaign per arm is missing.}
\label{tab:v3}
\begin{tabular}{lrrr}
\toprule
Arm & Fixed attack & Fixed utility & Adaptive attack\\
\midrule
Baseline & 7/28 & 11/28 & 11/27\\
Point-only & 4/28 & 14/28 & 4/27\\
Full history & 0/28 & 10/28 & 0/27\\
\bottomrule
\end{tabular}
\end{table}

Full history reduced fixed attack completion by 14.3 percentage points relative to point-only: four favorable discordances and none reversed. They span only two injection goals; the 20,000-replicate goal-cluster percentile sensitivity interval is $[-35.7,0.0]$ points. This supports a narrow history effect, not broad superiority.

The corrected adaptive workload produced 81/84 valid campaigns. A shared shard ledger reached 3,000,050 reported tokens, exceeding its three-million-token pre-call cap; the next calls stopped, leaving case v3-25 incomplete in all arms. These are missing outcomes, not blocked attacks. Among 27 complete pairs, full history prevented four point-only successes, with none reversed. Across all 28 planned pairs, worst-case missing-outcome bounds for full-history minus point-only attack risk are $[-17.9,-10.7]$ points. These finite-sample bounds are not population confidence intervals. The goal-cluster sensitivity interval is $[-37.0,0.0]$ points; the contrast still spans only two goals. Zero successes among valid full-history campaigns does not establish adaptive robustness.

The corrected clean-task replication completed all 40 workspace tasks per arm: baseline succeeded on 31, point-only on 30, and full history on 16. Full history lost 15 of 31 baseline-successful tasks, with no reverse gains, and denied 31 calls across 19/40 tasks. Relative to point-only, it lost 14 tasks with none gained: a $-35$-point utility difference (user-bootstrap sensitivity interval $[-50,-20]$). The separate initial clean run had the same aggregate success counts; we do not pool these repetitions. These observations expose a severe utility cost without assigning individual causality. No second model has been verified.

\section{Threats to Validity}
\textbf{Internal validity.} Language-model trajectories are stochastic. V2 arm order was fixed; v3 used balanced arm orders. Pairing controls task identity, not sampling variability. Shared task primitives and task heterogeneity limit generalization of nominal pair-level intervals and tests; no cluster-aware confirmatory inference was prespecified.

\textbf{Construct validity.} AgentDojo attack success represents completion of its defined injection objective. It is not equivalent to every form of compromise. Utility under attack depends on benchmark semantics and agent capability. Published work identifies AgentDojo injection-overwrite and evaluator limitations~\cite{bhagwatkar2025firewalls}; this is a validation concern, not proof that these cases were mis-scored.

\textbf{Policy validity.} Score weights and thresholds are engineered rules rather than calibrated probabilities. The confirmatory test evaluates the frozen rule set as a system, not the optimality of each weight. Post-hoc score replay quantifies decision changes but cannot resolve the agent behavior that would follow different interventions.

\textbf{Metadata validity.} Tool mappings, destination trust, sensitivity, provenance, and lineage are trusted inputs. Incorrect mappings can cause unsafe decisions or unnecessary blocking.

\textbf{External validity.} V2 used AgentDojo workspace, \texttt{tool\_knowledge}, one primary model deployment, and a fixed tool catalog. Results may not transfer to other suites, models, attacks, or production applications.

\textbf{Adaptive validity.} V2 did not expose the policy to an attacker that optimized prompts against its behavior; the corrected v3 study tests one bounded policy-aware attacker and retains three missing campaigns. It does not cover stronger attackers or independent attacker models. Published adaptive and automated attack work shows that this omission matters~\cite{zhan2025adaptive,hofer2026automated}.

\section{Conclusion}
Agent Flight Recorder enforces deterministic policy where model proposals become tool effects. In a protocol-frozen 120-pair AgentDojo workspace evaluation, the frozen system recorded 0/120 attack successes versus 20/120 for the paired baseline. The same policy reduced legitimate-task utility from 61/120 to 35/120. The evidence therefore shows both protection and cost.

Observable runtime state can support meaningful intervention without access to private model reasoning. The experimental result is narrower. It applies to the recorded benchmark and policy conditions, and it leaves adaptive robustness, metadata assurance, and usable human review as open engineering problems.

\appendices
\section{Ethical Considerations}
The study used synthetic benchmark environments, not production services or real users. Credentials were excluded from artifacts. The paper reports benchmark attacks and defense limits for reproducibility, without claiming safety beyond the measured conditions.

\section{Reproducibility}
Source, protocols, raw-evidence digests, and analyses are available at \url{https://github.com/upramod/agent-flight-recorder}. V2 is run 35180860582, commit \texttt{9ef3022}; all 240 endpoints match its archive. V3 separately evaluates repaired argument normalization and provenance propagation. The literal-payload replication is run 37571375916, commit \texttt{9bbcf7f}. Historical cells are retained, never pooled with corrected runs. Deterministic diagnostics need no credentials; live reproduction requires model access.

\section*{Acknowledgment}
OpenAI ChatGPT assisted with drafting and revising text throughout the manuscript, technical wording, repository and CI scripting, and presentation checks. Benchmark observations came from recorded experimental workflows; AI assistance did not supply experimental observations. The author is responsible for the manuscript and its claims.

\bibliographystyle{IEEEtran}
\bibliography{references}

\end{document}
